\documentclass{article}
\usepackage{graphicx}
\usepackage{amsmath}
\usepackage{float}
\usepackage{siunitx}
\usepackage{listings}
\usepackage{hyperref}
\usepackage{booktabs}
\usepackage{amsfonts}
\usepackage{textcomp}
\lstset{
    breaklines=true,
    breakatwhitespace=true,
    columns=fullflexible
}

\DeclareSIUnit[qualifier-mode = combine]{\dBm}{\deci\bel\of{m}}
\sisetup{exponent-product = \cdot, output-product = \cdot}

\title{IoT Challenge 3: Node-RED and LoRaWAN - Exercise}
\author{Angelica Ortiz 11191225, Daniel Nägele 11180610}
\date{April 2026}

\begin{document}

\maketitle
\thispagestyle{empty}
\newpage

\tableofcontents
\newpage


\section{EQ1}

The LoRaWAN network operates at 868 MHz with bandwidth 125 kHz. The network consists of 40 sensor nodes, each transmitting 2 packets per minute. Therefore, the total traffic load is:

\begin{equation}
\lambda_{tot} = 40 \cdot 2 = 80 \text{ packets/min}
\end{equation}

Converting to packets per second:

\begin{equation}
\lambda_{tot} = \frac{80}{60} = 1.333 \text{ packets/s}
\end{equation}

Using the TTN airtime calculator with payload size 20 bytes, the following airtime values were obtained:

\begin{center}
\begin{tabular}{cccc}
\toprule
SF & Airtime (s) & $G = \lambda T$ & $P_{success}$ \\
\midrule
SF7  & 0.0719 & 0.0959 & 82.55\% \\
SF8  & 0.1336 & 0.1781 & 70.03\% \\
SF9  & 0.2468 & 0.3291 & 51.78\% \\
SF10 & 0.4526 & 0.6035 & 29.91\% \\
SF11 & 0.9871 & 1.3161 & 7.19\% \\
SF12 & 1.8104 & 2.4139 & 0.80\% \\
\bottomrule
\end{tabular}
\end{center}

Assuming uncoordinated transmissions and a pure ALOHA model, the success probability is given by:

\begin{equation}
P_{success} = e^{-2\lambda T_{air}}
\end{equation}

From the results, only SF7 provides a success probability greater than 75\%. Therefore, the largest spreading factor that satisfies the requirement is:

\begin{equation}
\boxed{SF7}
\end{equation}

\section{EQ2}

The most appropriate action is: Move the nodes closer to the gateway

The non-uniform performance indicates that some nodes suffer from poor link quality due to distance or obstacles. Moving nodes closer to the gateway improves signal strength and reliability.

Increasing the spreading factor would increase airtime and collision probability, while decreasing the number of nodes does not solve the uneven performance across nodes.

\end{document}
